\questao{113}{corC}{Ondas e óptica}{Som em cordas: frequência e comprimento}{Com conta}
\begin{enunciado}
% gerado por ferramentas/md2tex_enunciado.py a partir de questoes/q113.md (não editar à mão)
O comprimento da parte oscilante da corda de um violão determina a nota musical emitida. No braço do violão, são colocadas peças metálicas espaçadas chamadas trastes. Quando as cordas são pressionadas entre os trastes, os comprimentos de suas partes oscilantes são reduzidos, gerando ondas de diferentes frequências. Algumas notas com os valores aproximados de suas frequências são apresentadas no quadro. A figura ilustra um violão com cordas de comprimento total vibrante ($L_T$) igual a 60\,cm. Considere o modo de vibração fundamental ($n = 1$) e que as frequências de vibração são inversamente proporcionais ao comprimento da parte oscilante da corda. Sabe-se que a primeira corda oscila com frequência de 660\,Hz (a nota mi) quando não pressionada (corda livre).\par
\begin{figlado}{0.4}{q113_fig1.png}
\textbf{Frequência das notas musicais, em hertz, no sistema temperado}\par
\begin{center}\renewcommand{\arraystretch}{1.15}
\begin{tabular}{cccccccc}
\toprule
\textbf{Oitava} & \textbf{Dó} & \textbf{Ré} & \textbf{Mi} & \textbf{Fá} & \textbf{Sol} & \textbf{Lá} & \textbf{Si}\\
\midrule
3ª & 130 & 147 & 165 & 175 & 196 & 220 & 248\\
4ª & 260 & 294 & 330 & 349 & 392 & 440 & 495\\
5ª & 520 & 587 & 660 & 698 & 784 & 880 & 990\\
6ª & 1\,040 & 1\,174 & 1\,320 & 1\,397 & 1\,568 & 1\,760 & 1\,980\\
\bottomrule
\end{tabular}\end{center}
\end{figlado}
\fonte{OLIVEIRA, Leonardo M.; VEIT, Eliane A.; SCHNEIDER, Claúdio. \textbf{Oitavas e notas}. Disponível em: www.if.ufrgs.br. Acesso em: 6 nov. 2023 (adaptado).}
Ao pressionar a primeira corda, de forma que a sua parte oscilante tenha 45 centímetros, e fazê-la vibrar, qual é a nota emitida pela corda?\par
\begin{altlinha}
\alth{A}{Dó.}%
\alth{B}{Mi.}%
\alth{C}{Sol.}%
\alth{D}{Lá.}%
\alth{E}{Si.}
\end{altlinha}

\end{enunciado}

\quer{qual nota a corda emite quando a parte que vibra é encurtada de 60\un{cm} para 45\un{cm}.}

\begin{ideia}
\textbf{Frequência} é o número de vibrações por segundo, medida em hertz (Hz). Frequência alta dá som agudo (fino);
frequência baixa dá som grave (grosso).

O enunciado entrega a regra: a frequência é \textbf{inversamente proporcional} ao comprimento da parte que vibra.
Ou seja: \textbf{corda mais curta, som mais agudo}.

\textbf{Faça o teste com uma régua.} Segure uma ponta da régua na beira da mesa e faça a outra ponta vibrar. Quanto menor o
pedaço para fora da mesa, mais rápido ele vibra e mais agudo é o som. É isso que o violonista faz ao apertar a corda.
(A régua serve só para sentir a ideia. A conta exata vale para a corda, como o enunciado garante.)

Quando duas medidas são inversamente proporcionais, \textbf{multiplicar uma pela outra dá sempre o mesmo resultado}:
\[ f \times L = \text{constante} \qquad\Rightarrow\qquad f_1 \times L_1 = f_2 \times L_2 \]

\begin{center}
\begin{tikzpicture}[font=\small]
  \draw[line width=1.4pt,corC] (0,1.1) -- (9,1.1);
  \fill (0,1.1) circle (2pt) (9,1.1) circle (2pt);
  \draw[cota] (0,0.75) -- (9,0.75) node[midway,below]{corda solta: 60\un{cm}\ \ $\rightarrow$\ \ 660\un{Hz}};
  \draw[line width=1.4pt,corB] (0,-0.5) -- (6.75,-0.5);
  \draw[line width=0.6pt,cinza,dashed] (6.75,-0.5) -- (9,-0.5);
  \fill (0,-0.5) circle (2pt) (6.75,-0.5) circle (2pt);
  \node[above,cinza,font=\scriptsize] at (6.75,-0.45) {dedo};
  \draw[cota] (0,-0.85) -- (6.75,-0.85) node[midway,below]{corda apertada: 45\un{cm}\ \ $\rightarrow$\ \ ?};
\end{tikzpicture}
\end{center}
\end{ideia}

\passoapasso
\npasso{Organize os dados.}
\begin{itemize}
  \item Corda solta: $L_1 = 60\un{cm}$ e $f_1 = 660\un{Hz}$ (nota mi).
  \item Corda apertada: $L_2 = 45\un{cm}$ e $f_2 =\ ?$
\end{itemize}
(Não se assuste com ``modo de vibração fundamental ($n = 1$)'' nem com ``sistema temperado'': esses termos não entram na conta.)

\npasso{Monte a igualdade: frequência $\times$ comprimento, antes e depois.}
\begin{conta}
$f_1 \times L_1 = f_2 \times L_2 \qquad\Rightarrow\qquad 660 \times 60 = f_2 \times 45$
\end{conta}

\npasso{Resolva.}
Queremos $f_2$ sozinho. Ele está multiplicado por 45; para desfazer, divida os dois lados por 45:
\begin{conta}
$f_2 = \dfrac{660 \times 60}{45}$
\end{conta}
Não multiplique ainda: simplifique primeiro. $\dfrac{60}{45} = \dfrac{4}{3}$ (dividindo os dois por 15: $60 \div 15 = 4$ e $45 \div 15 = 3$). Então:
\begin{conta}
$f_2 = 660 \times \dfrac{4}{3} = \dfrac{660}{3} \times 4 = 220 \times 4 = 880\un{Hz}$
\end{conta}
(Para fazer $660 \div 3$: $600 \div 3 = 200$ e $60 \div 3 = 20$. Somando: 220.)
\confere{a corda ficou mais curta, então a frequência tinha de \textbf{aumentar}. E aumentou: de 660 para 880\un{Hz}.}

\npasso{Procure 880\un{Hz} no quadro.}
Na linha da 5ª oitava, a coluna com 880 é a do \textbf{Lá}.

\resposta{D}{Com 45\un{cm}, a corda vibra a 880\un{Hz}, que é a nota Lá.}

\begin{outras}
\alt{B} Mi (660\un{Hz}) é a nota da corda solta. Ao encurtar a corda, a nota muda.\par
\alt{E} Si (495\un{Hz}) é o resultado de quem faz a proporção ao contrário (veja ``Onde muita gente erra'').\par
\altl{A} e \alt{C} As frequências de Dó e Sol não saem de nenhuma conta com esses dados.
\end{outras}

\begin{pegadinha}[Onde muita gente erra]
Fazer \textbf{regra de três direta}: $660 \times \dfrac{45}{60} = 495\un{Hz}$, que é um Si do quadro, e marcar E.
O erro está em esquecer a palavra \emph{inversamente}. Use o bom senso para conferir: corda \textbf{menor} tem de dar
frequência \textbf{maior}. Se a sua resposta deu menor que 660\un{Hz}, a proporção foi montada ao contrário.
\end{pegadinha}
